\documentclass[11pt]{article}

\usepackage[margin=1in]{geometry}
% T5 is the Vietnamese font encoding: the author line carries diacritics
% (ễ, ố, ắ) that pdflatex rejects under the default encoding, which is what
% made this file fail to build rather than anything about the machine. T5 is
% listed last so it stays the active encoding; T1 remains available for the
% rest of the document.
\usepackage[T1,T5]{fontenc}
\usepackage{amsmath,amssymb}
\usepackage{booktabs}
\usepackage{xcolor}
\usepackage[colorlinks=true,linkcolor=blue,citecolor=blue,urlcolor=blue]{hyperref}

\title{Comment on ``Descriptor: Heart and Lung Sounds Dataset Recorded from\\
a Clinical Manikin using Digital Stethoscope (HLS-CMDS)''}
\author{%
  Satya Adhiyaksa\\
  \small National Taiwan University of Science and Technology (NTUST), Taiwan\\
  \small \texttt{m11302825@mail.ntust.edu.tw}
  \and
  Nguyễn Quốc Thắng\\
  \small Ho Chi Minh City University of Technology (HCMUT),\\
  \small Vietnam National University Ho Chi Minh City, Vietnam
  \and
  Jenq-Shiou Leu\\
  \small National Taiwan University of Science and Technology (NTUST), Taiwan}
\date{\today}

\begin{document}
\maketitle

\begin{center}
\fbox{\parbox{0.9\textwidth}{\small
\textbf{Status: DRAFT, ready for review.} Both open items are now closed.
Provenance: the copy audited here is byte-identical to the Mendeley release
--- all six top-level artifacts match on MD5, verified 2026-08-19 --- so the
findings concern the released dataset rather than a redistribution. Sample
rate: checked against the descriptor's own text (Table~I, p.~138, and the
comparative claim on p.~134), which
changed how Section~\ref{sec:samplerate} is phrased --- 22\,050~Hz is stated
there as an instrument specification, and the descriptor never states the
released files' rate at all.}}
\end{center}

\begin{abstract}
While using the HLS-CMDS dataset \cite{hlscmds} to build an evaluation
pipeline for heart/lung sound source separation, we ran independent
integrity checks against the paired \texttt{Mix.csv} recordings that
\texttt{Mix.csv} claims to be additive combinations of the isolated
\texttt{HS.csv}/\texttt{LS.csv} recordings. Two independent methods --- a
least-squares additivity fit and a clipping-artifact scan --- agree that
only 36 of the 145 published mixture rows are acoustically consistent with
their stated heart/lung sources; the remaining 109 are not. We also report a
sample-rate gap: every released file measures 4{,}000~Hz, while the only
sample rate the descriptor gives is 22{,}050~Hz, stated as a specification
of the recording instrument; the section describing the released files does
not state their rate at all. We report both findings here, separately from
our own separation-methods work, because the first bears on ground truth
already in use downstream. Of the four published studies we are aware of
that build on HLS-CMDS, two evaluate on all 145 rows, and for one of them
the exposure is measurable from its own tables: 4{,}033 of the 5{,}365
scored windows behind its headline result fall on rows we find
non-additive (Section~\ref{sec:impact}).
\end{abstract}

\section{Motivation for this comment}

This comment is deliberately kept separate from our companion methods paper (in preparation)
on leakage-safe heart/lung separation evaluation. That paper's contribution
is a separation pipeline and baseline comparison; the findings below are
about the integrity of the released dataset itself, and are relevant to
every downstream user of HLS-CMDS, not only to our own results. We report
them here, addressed to the descriptor's own venue, so they are discoverable
by anyone consulting the dataset in the future.

\section{Additive-triplet audit}
\label{sec:additivity}

HLS-CMDS's \texttt{Mix.csv} lists, for each of 145 mixture recordings, the
isolated heart (\texttt{HS.csv}) and lung (\texttt{LS.csv}) recording IDs
that were combined to produce it. We tested this claim directly: for each
row, we fit the least-squares scalar gain $a$ minimizing
$\lVert \text{mixed} - a\cdot(\text{heart}+\text{lung}) \rVert$, then
computed the relative residual
$\lVert \text{mixed} - a\cdot(\text{heart}+\text{lung}) \rVert / \lVert \text{mixed} \rVert$.
A row is additive when this residual falls below $10^{-3}$.

The result is sharply bimodal, not a continuum: genuinely additive rows land
at relative residual $\sim 10^{-8}$--$10^{-4}$ (consistent with 16-bit
quantization noise on an exact sum), while the remaining rows land at
relative residual $\sim 1$ --- i.e., the mixed recording shares essentially
no structure with its named heart/lung sources. Only \textbf{36 of 145 rows}
(\texttt{M0087}, \texttt{M0111}--\texttt{M0145}) pass; the other 109 fail by
a wide margin, not a borderline one.

This check is implemented as \texttt{verify\_additive\_triplets()} in our
public repository (Section~\ref{sec:availability}) so it is independently
re-runnable against any copy of the dataset, rather than resting on a
one-off finding.

\section{Independent corroboration: a clipping fingerprint}
\label{sec:clipping}

A separate, unrelated audio-quality scan --- flagging the fraction of
samples pinned at the 16-bit full-scale rail ($\pm 32767$) in every
\texttt{.wav} file referenced by \texttt{HS.csv}/\texttt{LS.csv}/\texttt{Mix.csv}
--- turned up a single-sample clipping artifact in exactly one subset of
\texttt{Mix.csv}: \texttt{M0111}--\texttt{M0145} plus \texttt{M0087}. This is
the identical 36-row set identified by the additivity test in
Section~\ref{sec:additivity}, found independently, by a method with no
knowledge of the additivity result.

The mechanism is explicable and not an artifact of file corruption: it is
the fingerprint of peak normalization applied \emph{after} summing the
heart and lung signals, with a measured gain range of roughly
$3.5\times$--$135\times$ across these 36 rows (not integer overflow, which
was our original hypothesis before checking the gain magnitudes directly).
That two independent checks --- one purely numerical, one a physical
artifact of the recording pipeline --- converge on the same 36 IDs is, in
our view, strong evidence that this is a genuine property of the released
files rather than a bug in either check.

\section{Sample-rate discrepancy}
\label{sec:samplerate}

Every one of the 535 released files measures 4{,}000~Hz --- mono, 16-bit,
15.0~s (60{,}000 frames) --- without exception. The descriptor gives one
sample-rate figure, and it is $5.5\times$ higher: Table~I on page~138,
\emph{Technical Specification of the Digital Stethoscope}, lists
``Recording Sample Rate: 22\,050~Hz''.

We want to be precise about what that does and does not assert, because the
distinction changes what we are asking for. Table~I specifies the
\emph{instrument}, not the distributed files. The section that does describe
the distributed files --- \emph{Records and Storage} --- states their count,
format, grouping and metadata, but never their sample rate. The only other
passage that touches the rate is comparative: on page~134, immediately after
noting that the earlier manikin dataset of Tsai et al.\ was recorded ``at
8~kHz sampling rate'', the descriptor states that its own digital
stethoscope provides ``improved recording quality, a higher sampling rate,
and better resolution''. The released files, at 4{,}000~Hz, are below that
8~kHz reference, so this is the one sentence in the descriptor that the
release contradicts outright. Everything else is consistent with the files
having been captured at 22\,050~Hz and downsampled afterwards, undocumented.

What the descriptor leaves a reader with, then, is one explicit sample-rate
number and one comparative claim, and no way to learn that the files carry a
rate that satisfies neither. The recording path the
descriptor itself describes --- stethoscope, then Bluetooth to the Eko
mobile application, then that vendor's cloud platform, then download as
\texttt{.wav} --- has several points at which a resampling could occur, and
none of them is documented. A reader who takes 22\,050~Hz as the data rate,
which is the only inference the paper supports, will be wrong by a factor of
$5.5$.

We note explicitly that this is a property of the release and not of a
redistribution, because our own first reading of it was the opposite. The
copy we initially audited was obtained from
\url{github.com/Torabiy/HLS-CMDS}, which left open the likelier-seeming
explanation that we had measured a downsampled convenience mirror. We tested
that by re-downloading the dataset from Mendeley
(\url{https://data.mendeley.com/datasets/8972jxbpmp/3}, v3, retrieved
2026-08-19) and checksumming it independently: all six top-level artifacts
(\texttt{HS}, \texttt{LS} and \texttt{Mix}, each as \texttt{.csv} and
\texttt{.zip}) are byte-identical between the two copies. The mirror is
faithful; the discrepancy is in the release.

We therefore recommend the narrower fix that this finding actually supports:
that \emph{Records and Storage} state the released files' own sample rate
explicitly, and that whatever processing sits between capture at 22\,050~Hz
and release at 4\,000~Hz be described, or the files re-issued at the capture
rate if the downsampling was not intended.

The discrepancy is unlikely to have invalidated published results built on
this dataset --- cardiac energy below 200~Hz and respiratory energy below
1~kHz both sit comfortably under a 2~kHz Nyquist, so the retained band is
adequate for the analyses these papers perform. It does affect tooling and
reproduction: pretrained separation and classification models expecting
8--16~kHz input cannot be applied to these files without resampling, a
reader sizing an STFT or a filter bank from 22\,050~Hz will get every
frequency bin wrong, and anyone reproducing a published pipeline needs to
know which rate its author actually worked at.

\section{Downstream impact}
\label{sec:impact}

We are aware of four published studies that build on HLS-CMDS
\cite{aidriven,spectrotemporal,nmfcrnn,edgelung}. We determined, from each
paper's own text, tables and figures, which \texttt{Mix.csv} rows it
evaluates on (Table~\ref{tab:downstream}).

\begin{table}[htbp]
\centering
\small
\caption{\texttt{Mix.csv} rows entering each downstream study's evaluation,
as determined from that study's own text and tables. ``Outside the 36''
counts rows that fail the additivity test of
Section~\ref{sec:additivity}.}
\label{tab:downstream}
\begin{tabular}{@{}p{0.16\textwidth}p{0.40\textwidth}p{0.34\textwidth}@{}}
\toprule
Study & \texttt{Mix.csv} rows evaluated on & Outside the 36 additive rows \\
\midrule
Torabi \cite{aidriven} &
None. The separation chapters build their own mixtures: \S4.2.3 (p.~52)
re-pairs HLS-CMDS heart and lung segments at random into 25{,}000
mixtures; \S4.1.3 (p.~41) reports ``210 clinical manikin recordings'',
which is neither 145 nor identified as \texttt{Mix.csv}. \S5.5 (p.~65) and
\S6.2.2 (p.~72) use the 50 isolated heart and 50 isolated lung recordings
only. &
0. The \S4.1.3 evaluation set is \emph{not stated}; we record that as a
finding rather than infer it. \\
\addlinespace
Yaqub et al.\ \cite{spectrotemporal} &
All 145. \emph{Not stated in the text.} Table~9 (p.~2528) reports 5{,}365
scored windows on bandpass-separated heart sounds; segmentation is into
400~ms windows (p.~2514), giving 37 whole windows per 15~s recording, and
$145 \times 37 = 5365$. Its four per-class supports (481 / 2109 / 1147 /
1628) equal 37 times \texttt{Mix.csv}'s own per-class row counts (13 / 57 /
31 / 44) exactly, for every class. &
109 of 145 rows, i.e.\ \textbf{4{,}033 of the 5{,}365 windows (75.2\%)}, by
class: 333 Normal, 1{,}628 Murmur, 851 Extra Sound, 1{,}221 Rhythm
Disorder. \\
\addlinespace
Han et al.\ \cite{nmfcrnn} &
None scored. HLS-CMDS is ``used exclusively for dictionary learning and
[was] not included in the respiratory disease classification experiments''
(\S3.2.3, p.~7). \emph{Which} isolated recordings the dictionaries were fit
on is not stated. Figure~3 (p.~8) does show one mixture beside a
ground-truth respiratory signal; its provenance is not stated. &
0 scored; 1 unidentified row in a qualitative figure. \\
\addlinespace
Puneet et al.\ \cite{edgelung} &
All 145, \emph{stated}: the lung-isolation algorithm ``is evaluated on the
HLS CMS dataset,'' enumerated as Normal~(28), Wheezing~(28), Pleural
Rub~(25), Rhonchi~(23), Fine Crackles~(22), Coarse Crackles~(19) (p.~23) --
summing to 145 and matching \texttt{Mix.csv}'s lung-type distribution on
every class. &
109 of 145 rows (75.2\%): 18 Normal, 16 Wheezing, 24 Pleural Rub, 21
Rhonchi, 13 Fine Crackles, 17 Coarse Crackles. No separation metric is
reported on them; the 80.55\% accuracy (p.~24) is measured on \texttt{HF\_Lung}~V1
(p.~23), not on HLS-CMDS. \\
\bottomrule
\end{tabular}
\end{table}

Two of the four therefore evaluate on the full 145 rows, and so on 109 rows
whose mixed recording is not an additive combination of its listed sources.
For Yaqub et al.\ \cite{spectrotemporal} this is measurable: 4{,}033 of the
5{,}365 windows behind the 41.0\% accuracy of their Table~9 carry a heart
sound class label taken from a heart recording that our audit finds absent
from the audio being classified. Puneet et al.\ \cite{edgelung} run their
EVMD separation over the same 109 rows but report no metric against the
listed sources, so the exposure there is to the illustrated separation
quality rather than to a published number. The remaining two are not
exposed to the additivity finding: Torabi \cite{aidriven} constructs its own
mixtures rather than using \texttt{Mix.csv}'s, and Han et al.\
\cite{nmfcrnn} use HLS-CMDS only to fit dictionaries. None of the four
papers reports an independent integrity check of the kind described here,
and two of them -- Torabi's \S4.1.3 and Han et al.'s Figure~3 -- do not
state which recordings they used, so the exposure of those two cannot be
settled from the published text at all.

We raise this not to characterize any of these results as incorrect, but
because none of the four papers had the means to know whether it applied --
the additivity gap is not visible from the CSV metadata or from listening
to a handful of files; it required a systematic per-row check.

\section{Recommendation}

We suggest the descriptor be amended, or an erratum issued, to note: (1)
that only 36 of 145 \texttt{Mix.csv} rows are additive combinations of their
listed sources, with row IDs enumerated above, so that downstream users can
restrict evaluation to the valid subset or account for the rest --- two
published studies already evaluate on all 145
(Section~\ref{sec:impact}); and (2)
the sample-rate finding of Section~\ref{sec:samplerate} once resolved. We
would welcome the opportunity to share our audit code directly with the
authors if useful for reproducing or disputing this finding.

\section{Data and code availability}
\label{sec:availability}

The additivity check (\texttt{verify\_additive\_triplets()}) and the
clipping-based audio-quality scan referenced above are available at
\url{https://github.com/SouthsideProgramer/TEEP_Internship} (commit
\texttt{88c0ef8}), and are independently re-runnable against any copy of
HLS-CMDS: \texttt{make install}, \texttt{make dataset} (which verifies the
six release checksums), then the audit targets. No part of the dataset is
redistributed there.

\begin{thebibliography}{9}

\bibitem{hlscmds}
Y.~Torabi, S.~Shirani, and J.~P.~Reilly, ``Descriptor: Heart and Lung Sounds
Dataset Recorded from a Clinical Manikin using Digital Stethoscope
(HLS-CMDS),'' \textit{IEEE Data Descriptions}, vol.~2, pp.~133--140, 2025.
doi:10.1109/IEEEDATA.2025.3566012.

\bibitem{aidriven}
Y.~Torabi, ``AI-Driven Cardiorespiratory Signal Processing: Separation,
Clustering, and Anomaly Detection,'' Ph.D.\ dissertation, McMaster
University, 2025. Page numbers cited above refer to
\texttt{arXiv:2602.09210v1} (xxi\,+\,123~pp.), the author's posting of the
same thesis.

\bibitem{spectrotemporal}
A.~Yaqub, M.~S.~Orakzai, M.~F.~Qureshi, Z.~Mushtaq, I.~Siddique, and
T.~Radwan, ``Spectrotemporal Deep Learning for Heart Sound Classification
under Clinical Noise Conditions,'' \textit{Computer Modeling in
Engineering \& Sciences}, vol.~145, no.~2, pp.~2503--2529, 2025.
doi:10.32604/cmes.2025.071571.

\bibitem{nmfcrnn}
B.~Han, W.~Quan, B.~Matuszewski, and D.~Corbett, ``Respiratory Disease
Classification Using NMF-Enhanced Log-Mel Spectrograms and Convolutional
Recurrent Neural Networks,'' \textit{Sensors}, vol.~26, no.~13, art.~4268,
2026. doi:10.3390/s26134268.

\bibitem{edgelung}
A.~Puneet, P.~Shankar, S.~R.~R.~Koluguri, and A.~Srivastava,
``Edge-Enabled Portable Classifier for Lung Sounds Using Convolutional
Neural Networks,'' in \textit{Proc.\ 2025 IEEE Biomedical Circuits and
Systems Conference (BioCAS)}, pp.~21--25, 2025.
doi:10.1109/BioCAS67066.2025.00016.

\end{thebibliography}

\end{document}
